% =====================================================================
% DRAFT v1 (11.11.2026) -- compound border moves for the job shop
% Target: journal (C&OR / J. Scheduling / EJOR); plan B: short theory note.
% Layout: plain article class for drafting (as articles/kbonacci/main.tex).
% Plan: articles/jobshop/PLAN.md (v3, section A.7 = current state).
% ALL NUMBERS IN TABLES 2-4 ARE PILOT VALUES (1 s, Python prototype,
% 2-3 seeds). They show direction only and will be replaced by the
% campaign (stages min, opt, continuum, bridge).
% References: bibliographic data typed from memory, every entry must be
% checked against Crossref before submission.
% =====================================================================
\documentclass[10pt,twocolumn]{article}
\usepackage[margin=0.75in]{geometry}
\usepackage[T1]{fontenc}
\usepackage{amsmath,amssymb,amsthm}
\usepackage{booktabs}
\usepackage{multirow}
\usepackage{xcolor}
\usepackage[hidelinks]{hyperref}

\newtheorem{theorem}{Theorem}
\newtheorem{proposition}{Proposition}
\newtheorem{lemma}{Lemma}
\newtheorem{corollary}{Corollary}
\newtheorem{definition}{Definition}
\newcommand{\todo}[1]{\textcolor{red}{[TODO: #1]}}
\newcommand{\pilot}{\textcolor{blue}{[pilot]}}
\newcommand{\JP}{\mathrm{JP}}
\newcommand{\JS}{\mathrm{JS}}
\newcommand{\Cmax}{C_{\max}}

\title{Compound Border Moves for the Job Shop:\\
Feasibility of Block Reorderings and the\\
Flow-Shop-to-Job-Shop Continuum}
\author{\todo{authors}}
\date{}

\begin{document}
\maketitle

\begin{abstract}
Compound neighborhoods apply many independent swaps in one move and pick
the best combination by dynamic programming. They work well for the
permutation flow shop, where every permutation is a feasible schedule. We
carry them over to the job shop with makespan criterion. The critical path
takes the role of the permutation and its blocks take the role of the
segments in which swaps are placed. Feasibility is the one new obstacle. We
prove an exact condition for it: a reordering of operations inside critical
blocks creates a cycle if and only if it reverses a pair of operations that
is joined by a path leaving the first through its job successor, entering
the second through its job predecessor and avoiding the rest of the block.
The condition implies that compound moves built from disjoint intervals
need only tests of their single components, that crossing swaps add
pairwise constraints, and that the flow shop has no infeasible reordering
at all. The classical block property then reduces every swap family to at
most two moves per block. This leads to a family of border moves with three
axes: move type (swap or insertion), reach, and single versus compound
application. N5, N6 and N7 are points of this family, and so are the
compound Fibonacci, $k$-bonacci and Dynasearch neighborhoods. In a pilot
study with iterated local search, compound moves of moderate reach attain a
lower makespan than N5 and N6 on Taillard instances, and their advantage
over N5 is largest for instances between the flow shop and the job shop.
\todo{replace the last sentence with the campaign result}
\end{abstract}

% ---------------------------------------------------------------------
\section{Introduction}

Local search for the job shop scheduling problem with makespan criterion,
$J||\Cmax$, rests on a small set of neighborhoods defined on the critical
path. N1 swaps any two adjacent critical operations of one
machine~\cite{laarhoven1992}. N5 keeps only the swaps at the two ends of
each critical block~\cite{nowicki1996}. N6 and N7 move an operation to the
front or to the back of its block~\cite{balas1998,zhang2007}. All of them
change one block per move, and the best tabu search
methods~\cite{nowicki2005,peng2015} are built on them.

For the permutation flow shop a different line of neighborhoods exists. A
compound move is a set of swaps on disjoint intervals of the permutation,
and the best set is found by dynamic programming, as in
Dynasearch~\cite{congram2002}. Restricting the swap distance gives the
Fibonacci neighborhood (adjacent swaps only)~\cite{fibonacci_own} and the
$k$-bonacci family between Fibonacci and
Dynasearch~\cite{kbonacci_own}. These neighborhoods are exponential in
size and cheap to search, but they rely on the fact that every permutation
is a feasible flow shop schedule. In the job shop a swap of two
non-adjacent operations may close a cycle in the disjunctive graph, and a
combination of feasible swaps is not obviously feasible.

This paper asks what remains of the compound neighborhoods when the routes
of the jobs differ. Our answer has three parts. The first is a
construction. The critical path is treated as a virtual permutation, a
candidate swap lives inside one block, and the selection of a compound
move is the same dynamic program as in the flow shop. The program
decomposes into the blocks, so the neighborhood is a product of block
neighborhoods. The second part is a feasibility theorem for arbitrary
reorderings inside blocks. It gives an exact, pairwise condition and
explains which compound families are safe: for moves on disjoint or nested
intervals it is enough to test every swap alone, crossing swaps need
pairwise tests, and with identical routes nothing needs testing. The third
part is an honest account of what the block property does to these
families. A move that changes no block border cannot improve the makespan,
so under the usual filter every block keeps at most two swaps, one at each
end. The sequence-counted families collapse, and the natural family for
the job shop is indexed by the type of the border move, by its reach and
by whether one move or one move per block end is applied. N5, N6 and N7
are the single-move points of this family. Their compound versions are
new.

The experimental part measures three things under one protocol and one
code path: the quality reached by each member of the family in iterated
local search and simulated annealing on the Taillard job shop
instances~\cite{taillard1993}, the share of infeasible long moves as a
function of instance shape and solution quality, and the behaviour of the
family on generated instances that interpolate between a flow shop and a
job shop. \todo{one sentence with the headline result of the campaign}

% ---------------------------------------------------------------------
\section{Background}

An instance has $n$ jobs and $m$ machines. Job $j$ visits the machines in
the order of its route and its $t$-th operation has processing time
$p_{j,t}$. For an operation $o$ we write $\JP(o)$ and $\JS(o)$ for its job
predecessor and successor. A solution is a sequence $\sigma_k$ of the
operations of every machine $k$. The disjunctive
graph~\cite{roy1964} $G(\sigma)$ has a job arc $o\to\JS(o)$ for every
operation and a machine arc between consecutive operations of every
$\sigma_k$. The solution is feasible if and only if $G(\sigma)$ is
acyclic. The head $r(o)$ is the length of a longest path ending at the
start of $o$, the tail $q(o)$ is the length of a longest path leaving the
end of $o$, and $\Cmax=\max_o\, r(o)+p_o+q(o)$. One topological pass
computes all of them in $O(nm)$ time.

A critical path is a longest path in $G(\sigma)$. A block is a maximal run
of consecutive critical operations joined by machine arcs. The block
property~\cite{brucker1994,nowicki1996} states that a schedule with a
smaller makespan must process some operation of a block before the first
or after the last operation of that block. N5 uses it to keep two swaps
per block. N6 and N7 use it to define insertions, together with sufficient
conditions for their feasibility~\cite{balas1998,zhang2007}. Feasibility
of insertions was studied in a general setting by Gr\"oflin and
Klinkert~\cite{groflin2007} and for regular criteria by Mati et
al.~\cite{mati2011}. \todo{state precisely how Theorem~\ref{th:main}
relates to \cite{groflin2007,mati2011,dauzere1997} and to the compound
moves of Bo\.zejko and Wodecki once these are read}
Candidate moves are usually ranked by an estimate computed from the
current heads and tails~\cite{taillard1994,balas1998} and only the chosen
move is evaluated exactly.

In the permutation flow shop the compound neighborhoods are counted by
integer sequences: a permutation of length $L$ has $F_{L+1}$ sets of
disjoint adjacent swaps, $a_k(L)$ sets of disjoint swaps of distance below
$k$ ($k$-bonacci numbers) and $2^{L-1}$ sets of disjoint swaps of any
distance. Allowing nested swaps gives the Motzkin numbers and allowing
crossings gives the involutions~\cite{motzkin_own}.

% ---------------------------------------------------------------------
\section{Neighborhoods on the Critical Path}\label{sec:family}

Let $P=(o_1,\dots,o_L)$ be one critical path of a feasible solution and
let $B_1,\dots,B_b$ be its blocks. When two predecessors of an operation
are critical we follow the machine predecessor, which gives longer blocks.
The positions $1,\dots,L$ of $P$ form a virtual permutation. A swap
exchanges two operations of one block and corresponds to an interval of
positions. Intervals of different blocks are disjoint, so any set of swaps
taken from different blocks is a compound move in the flow shop sense. A
flow shop move changes the order on all machines at once. A job shop move
changes one machine per block.

\begin{proposition}[size]\label{pr:count}
Let $a_F(L)$ be the number of compound moves of a family $F$ on a
permutation of length $L$. On a critical path with block lengths
$L_1,\dots,L_b$ the family has $\prod_{i} a_F(L_i)$ moves, including the
empty one.
\end{proposition}

The product is exponential in the path length, as in the flow shop. The
block property removes most of it.

\begin{proposition}[border rule]\label{pr:border}
Let $\sigma'$ be a feasible solution obtained from $\sigma$ by reordering
operations inside the blocks of $P$. If $\sigma'$ keeps the first
operation of every block except $B_1$ and the last operation of every
block except $B_b$, then $\Cmax(\sigma')\ge\Cmax(\sigma)$.
\end{proposition}
\begin{proof}
In $\sigma'$ the operations of each block are still consecutive on their
machine. Visit them in the new order. The job arc from the last operation
of $B_i$ to the first operation of $B_{i+1}$ is unchanged, so the visits
form a path in $G(\sigma')$ with the same operations as $P$ and the same
length.
\end{proof}

A swap that touches no border has a non-negative effect, so a dynamic
program that adds the effects of its swaps never selects it. We therefore
keep only candidates that replace the first or the last operation of a
block. This filter is exact for the additive model used in the selection.
It is not exact for the true makespan of a compound move, because an
interior swap can change the result of a border swap applied with it.
Section~\ref{sec:eval} returns to this point.

\begin{proposition}[flattening]\label{pr:flat}
Under the border rule a block of length $L$ contributes at most two
disjoint swaps to a compound move, one containing its first and one
containing its last position. Nested swaps cannot occur, and two swaps
cross only as the pair $(1,j)$, $(i,L)$ with $i<j$.
\end{proposition}

Proposition~\ref{pr:flat} changes the picture inherited from the flow
shop. Inside a block the Motzkin family equals Dynasearch, and the number
of filtered compound moves of a block is polynomial in $L$. What is left
to vary is the border move itself. We consider three types of moves on a
block interval $[a,b]$ that contains a block end: the endpoint swap, the
insertion that brings an inner operation to the block end (N6), and the
insertion that sends the end operation inward (N7). The reach parameter
$k$ bounds the interval by $b-a\le k-1$. For $k=2$ all three types are the
same adjacent swap. A single neighborhood applies the best candidate. A
compound neighborhood applies at most one candidate per block end in every
block at once. Table~\ref{tab:family} lists the members studied here.

\begin{table}[t]
\centering\small
\begin{tabular}{@{}lll@{}}
\toprule
Type and reach & Single & Compound \\
\midrule
swap, $k=2$ & N5 & Fib \\
swap, reach $k$ & & $k$-bonacci \\
swap, any reach & & Dyn \\
N6 insertion, any reach & N6 & c6 \\
N6 and N7 insertions, any reach & N7 & c7 \\
all types, reach $k$ & & bm$_k$ \\
\bottomrule
\end{tabular}
\caption{The border-move family. Entries c6, c7 and bm$_k$ are new.}
\label{tab:family}
\end{table}

One iteration evaluates the current solution, extracts the path and its
blocks, lists the border candidates of the chosen type and reach, discards
the infeasible ones (Section~\ref{sec:feas}), estimates the makespan after
each remaining candidate, selects a set of candidates and evaluates the
resulting solution exactly. The selection is the positional dynamic
program of the flow shop code. Because candidates of different blocks
never overlap, it returns the sum of independent block optima, and we do
not claim it as a contribution. If it selects nothing, the best single
candidate is applied, as in our flow shop studies. The cost of an
iteration is two $O(nm)$ evaluations plus $O(bk^2)$ for the estimates of a
family of reach $k$.

% ---------------------------------------------------------------------
\section{Feasibility of Block Reorderings}\label{sec:feas}

Let $\sigma$ be feasible and let $\sigma'$ be obtained by permuting the
operations inside blocks of $P$. Several blocks may change at once. For
two operations $x,y$ of one block, with $x$ before $y$ in $\sigma$, a
\emph{bypass} from $x$ to $y$ is a path in $G(\sigma)$ that leaves $x$
through the job arc $x\to\JS(x)$, enters $y$ through the job arc
$\JP(y)\to y$ and contains no other operation of the block. The case
$\JS(x)=\JP(y)$ is allowed.

\begin{lemma}\label{le:clean}
A bypass from $x$ to $y$ contains no operation of $P$ other than $x$ and
$y$.
\end{lemma}
\begin{proof}
Let $u$ be an inner operation of the bypass that lies on $P$. If $u$
precedes $x$ on $P$, then $P$ leads from $u$ to $x$ and the bypass leads
from $x$ to $u$, which is a cycle in $G(\sigma)$. If $u$ follows $y$ on
$P$, the cycle is formed by $P$ from $y$ to $u$ and the bypass from $u$ to
$y$. Otherwise $u$ lies between $x$ and $y$ on $P$, so it belongs to the
block, which the definition excludes.
\end{proof}

\begin{theorem}\label{th:main}
The solution $\sigma'$ is infeasible if and only if there are operations
$x,y$ in one block such that $x$ precedes $y$ in $\sigma$, $y$ precedes
$x$ in $\sigma'$, and $G(\sigma)$ contains a bypass from $x$ to $y$.
\end{theorem}
\begin{proof}
Write $\prec$ and $\prec'$ for the machine orders of $\sigma$ and
$\sigma'$. They differ only on pairs of operations of the same block.

\emph{Sufficiency.} By Lemma~\ref{le:clean} the inner arcs of the bypass
join operations outside $P$, and its first and last arcs are job arcs.
None of these arcs depends on the order inside a block, so the bypass is a
path of $G(\sigma')$. With $y\prec' x$ it closes a cycle.

\emph{Necessity.} Let $C$ be a cycle of $G(\sigma')$, where machine arcs
are taken in the transitive sense. Call a machine arc $u\to v$ of $C$
\emph{new} if $v\prec u$. The other arcs of $C$ belong to $G(\sigma)$, so
$C$ contains a new arc and splits into new arcs and paths of $G(\sigma)$
between them. A new arc stays inside one block and goes backwards on $P$.
A path of $G(\sigma)$ between two operations of $P$ cannot go backwards on
$P$. Once such a path left a block for a later one, $C$ could not return,
so all new arcs of $C$ lie in one block $B$.

Cut $C$ at every visit to $B$. A piece is either a single machine arc
between operations of $B$, ordered by $\prec'$, or a path from $u\in B$ to
$v\in B$ whose inner operations are outside $B$. A piece of the second
kind lies in $G(\sigma)$, so $u\prec v$. It cannot start with a machine
arc, because that arc would lead to an operation after the block on the
same machine and the rest of the piece would return to $B$, closing a
cycle in $G(\sigma)$. For the same reason it cannot end with a machine
arc. It is therefore a bypass from $u$ to $v$. Merging consecutive machine
pieces, $C$ reads
$c_1\preceq' d_1\leadsto c_2\preceq' d_2\leadsto\dots\leadsto c_1$, where
every $\leadsto$ is a bypass. At least one bypass occurs, since $\prec'$
is a linear order. If $d_i\prec' c_{i+1}$ held for every bypass, the
sequence would be a cycle of $\prec'$. Hence some bypass joins a pair with
$d_i\prec c_{i+1}$ and $c_{i+1}\prec' d_i$.
\end{proof}

The exclusion of the other block operations matters. If the path from $x$
to $y$ may pass through operations $z,z'$ of the block and use the machine
arc $z\to z'$, the new order $z',y,x,z$ reverses $(x,y)$ and is still
feasible. In brute-force tests on random instances the condition of
Theorem~\ref{th:main} agreed with cycle detection in all $184\,461$
reorderings tested, of which $28\,065$ were infeasible, while the version
without the exclusion gave one false alarm. \pilot

The theorem depends on $\sigma'$ only through its set of reversed pairs.
This gives the following consequences.

\begin{corollary}\label{co:all}
(a) A swap of two adjacent critical operations is
feasible~\cite{laarhoven1992}.
(b) A compound move whose components act on disjoint or nested intervals
is feasible if and only if every component is feasible alone.
(c) Moving $u$ forward past $x_1,\dots,x_s$ is infeasible if and only if
a bypass leads from $u$ to some $x_i$. Moving $u$ backward is infeasible
if and only if a bypass leads from some $x_i$ to $u$.
(d) If all jobs have the same route, every reordering inside blocks is
feasible.
\end{corollary}
\begin{proof}
(a) A bypass from $x$ to its machine successor $y$ is longer than
$p_x+p_y$, which contradicts the criticality of the arc $x\to y$. (b) The
set of reversed pairs of such a move is the union of the sets of its
components. (c) The move reverses exactly the pairs $(u,x_i)$ or
$(x_i,u)$. (d) With a common route the machine index never decreases
along a path, and a bypass first increases it.
\end{proof}

Part (b) covers the compound Fibonacci, $k$-bonacci, Dynasearch and
Motzkin families and all compound border moves of
Table~\ref{tab:family}: their feasibility costs one test per candidate
and nothing per combination. Crossing swaps behave differently. Their set
of reversed pairs is not the union of the components, two feasible swaps
may give an infeasible move (0.71\% of crossing pairs in our tests
\pilot), and an infeasible swap may be repaired by a crossing one. Every
reversed pair depends on at most two swaps, so feasibility of a set of
crossing swaps is a quadratic constraint in the selection variables.
Part (c) is exact, whereas the conditions used with N6 and
N7~\cite{balas1998,zhang2007} are sufficient. Part (d) identifies the
flow shop as the case without obstacles.

In the implementation a pair is first checked with heads: a path from
$\JS(x)$ to $\JP(y)$ requires $r(\JP(y))\ge r(\JS(x))+p_{\JS(x)}$. This
certificate settles about 90\% of the feasible non-adjacent swaps
\pilot. The remaining pairs are decided by a depth-first search from
$\JS(x)$ that is pruned by heads and never enters the block.

% ---------------------------------------------------------------------
\section{Move Evaluation and Additivity}\label{sec:eval}

For a candidate that reorders a block interval, the estimate recomputes
the heads of the interval operations from left to right and their tails
from right to left in the new order. Heads of job predecessors and of the
machine predecessor of the interval, and tails of job successors and of
the machine successor, are taken from the current solution. The estimate
of the new makespan is the largest value of $r'+p+q'$ in the interval and
costs time linear in the interval length. For an adjacent swap this is the
estimate of Taillard~\cite{taillard1994}: a lower bound on the new
makespan that is exact whenever the move does not improve. For longer
intervals it is an approximation. \todo{E0: rank correlation with the
exact value and share of false improvements, per reach}

The selection adds the estimated effects of the chosen candidates. In the
job shop this sum overstates the real gain more than in the flow shop,
because shortening one block makes a parallel path critical. In the pilot
the real gain of a compound move was 36--43\% of the sum of the exact
single effects and 14--19\% of the sum of the estimates, and still only
2--4\% of the compound moves failed to improve. \pilot{} The exact
evaluation of the selected move is therefore required. \todo{E5: remedies,
i.e. exact evaluation of the top-$t$ moves, a cap on the number of
components, and the unfiltered variant that measures what the border
filter loses}

% ---------------------------------------------------------------------
\section{Experimental Protocol}\label{sec:setup}

The protocol follows our flow shop studies so that the neighborhood is the
only factor that changes. Iterated local search and simulated annealing
keep their control flow and constants. \todo{constants in one sentence,
as in the $k$-bonacci paper} The start solution assigns the operations in
round-robin order of the jobs and decodes them semi-actively, which is the
analogue of the identity permutation. The perturbation is a double bridge
on a topological order of the operations followed by the same decoding, so
it always returns a feasible solution. Tabu identifiers are sets of
swapped operation pairs. All neighborhoods share one implementation that
differs only in the candidate generator.

The main benchmark is the set ta01--ta80 of Taillard~\cite{taillard1993},
eight classes from $15\times15$ to $100\times20$ with ten instances each.
We regenerated all instances from the published seeds. Each configuration
is run once for 10~s with five seeds, and the results for the budgets 0.1,
0.5, 1, 2 and 5~s are read from the move log of that run. We report the
relative deviation from the best known upper bound~\cite{vanhoorn2018}
with the date and source of every bound, and the number of iterations
next to every time budget. \todo{machine, Python version; validation of
the budget extraction on two classes}

The continuum experiment uses Taillard's generator with a route disorder
parameter $\rho\in\{0,0.1,0.25,0.5,1\}$. Routes start as the identity,
and each transposition of the route shuffle is performed with probability
$\rho$. Processing times do not depend on $\rho$. The value $\rho=0$ gives
a flow shop in which machine orders may differ, and $\rho=1$ gives the
standard job shop instance. Sizes are $20\times5$, $20\times10$,
$30\times15$ and $50\times10$ with ten instances each. As a bridge to the
flow shop literature, the Taillard flow shop instances of sizes
$20\times5$, $20\times10$ and $50\times10$ are also solved as job shops
with identical routes.

Variants are compared with the Friedman test and the Nemenyi critical
difference~\cite{demsar2006} over instances, and key pairs with the
Wilcoxon signed-rank test under the Holm correction at $\alpha=0.05$.

% ---------------------------------------------------------------------
\section{Results}\label{sec:exp}

\todo{This section holds pilot values only. Final structure: descent to
local optima; anatomy of the path and feasibility; main comparison on
ta01--ta80 with a critical difference diagram; continuum; additivity;
tabu search reference.}

\paragraph{Anatomy.}
Table~\ref{tab:anatomy} describes the critical path of solutions reached
by a short iterated local search. Blocks become longer as $n/m$ grows.
Few candidates survive the border rule for small reach, and a large share
of the long swaps is infeasible in the classes with $n/m\ge2$. The reach
that can be used is therefore bounded by feasibility and not only by the
cost of evaluation.

\begin{table}[t]
\centering\small
\begin{tabular}{@{}lrrrrr@{}}
\toprule
Class & blocks & length & $k{=}4$ & $k{=}\infty$ & infeas. \\
\midrule
$15\times15$  & 7.7  & 3.1  & 12.2 & 18.3  & 19\% \\
$20\times20$  & 13.0 & 2.4  & 18.7 & 28.2  & 20\% \\
$30\times15$  & 4.7  & 9.1  & 14.0 & 41.8  & 45\% \\
$50\times15$  & 3.7  & 20.1 & 11.7 & 63.7  & 57\% \\
$100\times20$ & 3.7  & 32.8 & 12.3 & 116.7 & 67\% \\
\bottomrule
\end{tabular}
\caption{Blocks, mean block length, swap candidates for reach 4 and
unbounded reach, and the infeasible share of the latter. \pilot}
\label{tab:anatomy}
\end{table}

\paragraph{Border-move family.}
Table~\ref{tab:ils} compares nine members of the family in iterated local
search. Compound swaps of reach 3 were better than N5 on all eleven
instances and better than N6 on nine. The compound version of N7 was
better than N7 on ten instances. For N6 the compound version won on seven
instances and lost on three, which a sign test does not separate. The
compound variants are close to each other. The pilot therefore suggests
that applying a border move in all blocks at once matters more than the
type of the move, and leaves the question open for N6.

\begin{table}[t]
\centering\small
\begin{tabular}{@{}lrrr@{}}
\toprule
Variant & RPD [\%] & mean rank & iter./s \\
\midrule
N5        & 10.62 & 8.77 & 2868 \\
N6        & 8.41  & 5.73 & 2304 \\
N7        & 8.90  & 7.05 & 1970 \\
$k=3$     & 7.40  & 3.36 & 2659 \\
Dyn       & 7.47  & 4.45 & 2095 \\
c6        & 7.53  & 3.68 & 2109 \\
c7        & 7.50  & 4.32 & 1974 \\
bm$_3$    & 7.42  & 3.68 & 2403 \\
bm$_\infty$ & 7.70 & 3.95 & 1592 \\
\bottomrule
\end{tabular}
\caption{Iterated local search, 1~s, eleven instances, three seeds;
deviation from the best known makespan. \pilot}
\label{tab:ils}
\end{table}

\paragraph{From flow shop to job shop.}
Table~\ref{tab:rho} shows the two larger sizes of the continuum
experiment. With identical routes N5 gave the best result for every size.
For $\rho$ between 0.1 and 0.5 it lost 14 to 22 percentage points to the
best variant, while the compound variants stayed within a few points of
each other. At $\rho=1$ the gap narrowed, and for $n/m\ge4$ it vanished
because the critical path of a good solution is a single block and the
makespan equals the machine bound. Two effects have to be separated
before this pattern is interpreted. N5 performs two to four times more
iterations than the compound variants at $\rho=0$, and mild route disorder
lengthens the blocks. \todo{continuum stage: results at equal iteration
counts, block lengths and infeasible shares per $\rho$}

\begin{table}[t]
\centering\small
\begin{tabular}{@{}lrrrrr@{}}
\toprule
Size & $\rho$ & N5 & $k{=}3$ & Dyn & c6 \\
\midrule
\multirow{5}{*}{$30\times15$}
 & 0    & 0.68  & 2.46 & 2.27 & 2.47 \\
 & 0.1  & 18.74 & 3.39 & 7.52 & 4.13 \\
 & 0.25 & 16.62 & 4.02 & 4.26 & 2.20 \\
 & 0.5  & 14.18 & 3.85 & 4.28 & 2.61 \\
 & 1    & 5.78  & 2.81 & 1.84 & 1.59 \\
\midrule
\multirow{5}{*}{$50\times10$}
 & 0    & 0.64  & 1.41 & 1.65 & 1.11 \\
 & 0.1  & 22.14 & 7.54 & 4.37 & 5.63 \\
 & 0.25 & 17.75 & 4.58 & 2.69 & 4.39 \\
 & 0.5  & 17.35 & 5.72 & 1.75 & 2.63 \\
 & 1    & 1.59  & 0.48 & 0.00 & 0.00 \\
\bottomrule
\end{tabular}
\caption{Deviation [\%] from the best result of any variant on the same
instance for route disorder $\rho$. \pilot}
\label{tab:rho}
\end{table}

% ---------------------------------------------------------------------
\section{Remark on Quantum Formulations}

The selection of a compound move is a quadratic unconstrained binary
problem with one variable per candidate and a penalty for every pair of
overlapping candidates, as in our flow shop work~\cite{windowed_own}. In
the job shop this problem is block diagonal and small: a block
contributes at most $2(k-1)$ variables and whole moves had 4 to 19
variables for $k\le4$ in good solutions. \pilot{} A whole move fits a
current gate-model device without windowing. The dynamic program solves
the same problem exactly, so no speed-up is claimed. The crossing
families are the interesting case, because their feasibility constraints
are true two-body terms. \todo{decide whether this section stays or moves
to a separate paper}

% ---------------------------------------------------------------------
\section{Conclusions}

\todo{Write after the campaign. Expected content: recommended reach as a
function of $n/m$; the role of compounding versus move type; where on the
flow-shop-to-job-shop axis compound moves pay off; limits, namely the
iterated local search framework, the Python implementation, and the
makespan criterion.}

\begin{thebibliography}{99}\small
\bibitem{balas1998} E.~Balas, A.~Vazacopoulos. Guided local search with
shifting bottleneck for job shop scheduling. \emph{Manag. Sci.}
44(2):262--275, 1998.
\bibitem{brucker1994} P.~Brucker, B.~Jurisch, B.~Sievers. A branch and
bound algorithm for the job-shop scheduling problem. \emph{Discrete Appl.
Math.} 49:107--127, 1994.
\bibitem{congram2002} R.K.~Congram, C.N.~Potts, S.L.~van de Velde. An
iterated dynasearch algorithm for the single-machine total weighted
tardiness scheduling problem. \emph{INFORMS J. Comput.} 14(1):52--67,
2002.
\bibitem{dauzere1997} S.~Dauz\`ere-P\'er\`es, J.~Paulli. An integrated
approach for modeling and solving the general multiprocessor job-shop
scheduling problem using tabu search. \emph{Ann. Oper. Res.} 70:281--306,
1997. \todo{verify}
\bibitem{demsar2006} J.~Dem\v{s}ar. Statistical comparisons of classifiers
over multiple data sets. \emph{J. Mach. Learn. Res.} 7:1--30, 2006.
\bibitem{fibonacci_own} \todo{own work: Fibonacci neighborhood, Sci. Rep. 2026}
\bibitem{groflin2007} H.~Gr\"oflin, A.~Klinkert. Feasible insertions in
job shop scheduling, short cycles and stable sets. \emph{Eur. J. Oper.
Res.} 177(2):763--785, 2007.
\bibitem{kbonacci_own} \todo{own work: $k$-bonacci family, under review}
\bibitem{laarhoven1992} P.J.M.~van Laarhoven, E.H.L.~Aarts, J.K.~Lenstra.
Job shop scheduling by simulated annealing. \emph{Oper. Res.}
40(1):113--125, 1992.
\bibitem{mati2011} Y.~Mati, S.~Dauz\`ere-P\'er\`es, C.~Lahlou. A general
approach for optimizing regular criteria in the job-shop scheduling
problem. \emph{Eur. J. Oper. Res.} 212(1):33--42, 2011.
\bibitem{motzkin_own} \todo{own work: Motzkin neighborhood, SOCO 2026}
\bibitem{nowicki1996} E.~Nowicki, C.~Smutnicki. A fast taboo search
algorithm for the job shop problem. \emph{Manag. Sci.} 42(6):797--813,
1996.
\bibitem{nowicki2005} E.~Nowicki, C.~Smutnicki. An advanced tabu search
algorithm for the job shop problem. \emph{J. Sched.} 8(2):145--159, 2005.
\bibitem{peng2015} B.~Peng, Z.~L\"u, T.C.E.~Cheng. A tabu search/path
relinking algorithm to solve the job shop scheduling problem.
\emph{Comput. Oper. Res.} 53:154--164, 2015.
\bibitem{roy1964} B.~Roy, B.~Sussmann. Les probl\`emes d'ordonnancement
avec contraintes disjonctives. Note DS no.~9 bis, SEMA, Paris, 1964.
\bibitem{taillard1993} E.~Taillard. Benchmarks for basic scheduling
problems. \emph{Eur. J. Oper. Res.} 64(2):278--285, 1993.
\bibitem{taillard1994} E.~Taillard. Parallel taboo search techniques for
the job shop scheduling problem. \emph{ORSA J. Comput.} 6(2):108--117,
1994.
\bibitem{vanhoorn2018} J.J.~van Hoorn. The current state of bounds on
benchmark instances of the job-shop scheduling problem. \emph{J. Sched.}
21(1):127--128, 2018.
\bibitem{windowed_own} \todo{own work: windowed QUBO neighborhoods}
\bibitem{zhang2007} C.Y.~Zhang, P.G.~Li, Z.L.~Guan, Y.Q.~Rao. A tabu
search algorithm with a new neighborhood structure for the job shop
scheduling problem. \emph{Comput. Oper. Res.} 34(11):3229--3242, 2007.
\end{thebibliography}

\end{document}
